% Commandes pour les flèches textuelles
% Désactive les shorthands dans TikZ
\AtBeginDocument{\shorthandoff{:;!?}}
\tikzset{every picture/.style={execute at begin picture={\shorthandoff{:;!?};}}}
\tikzstyle{every picture}+=[remember picture]
\tikzstyle{na} = [shape=rectangle,inner sep=0pt]


\newcommand{\ptFleche}[2]{		% Déclaration d'une extrémité de flèche
    \shorthandoff{:;!?}% <- désactive les shorthands
    \tikz[baseline=(#1.base)]\node[na](#1){#2};
  }

\newcommand{\Fleche}[5][thick]{%
    \begin{tikzpicture}[overlay]
        \path[->,#1](#2) edge [out=#4, in=#5] (#3);
    \end{tikzpicture}
}

% --- Commande FlecheEllipse ---
% #1 = nœud source (nom)
% #2 = angle sortie du nœud source
% #3 = angle entrée du nœud cible
% #4 = décalage dx, dy sous forme (dx,dy) en cm
% #5 = couleur ellipses et flèche
% #6 = texte du nœud cible B
\newcommand{\FlecheEllipseo}[6]{%
    \begin{tikzpicture}[overlay, remember picture]
        % Coordonnées source
        \path let \p1 = (#1) in 
        coordinate (Acoord) at (\x1,\y1);

        % Coordonnées cible
        \pgfmathsetmacro{\dx}{#4}
        \pgfmathsetmacro{\dy}{#5} % ici on sépare les deux pour le calcul

        \coordinate (B) at ([shift={(#4)}]#1);

        % Nœud A entouré d'ellipse
        \node[draw,#5,ellipse,inner sep=2pt] (#1-node) at (#1) {#1};
        % Nœud B entouré d'ellipse
        \node[draw,#5,ellipse,inner sep=2pt] (B-node) at (B) {#6};
        % Flèche de A vers B
        \draw[->,#5,thick] (#1-node) edge [out=#2,in=#3] (B-node);
    \end{tikzpicture}%
}

% --- Commande FlecheEllipse finale ---
% #1 = nom du nœud source (A)
% #2 = angle sortie du nœud source
% #3 = angle entrée du nœud cible
% #4 = décalage horizontal dx (en cm)
% #5 = décalage vertical dy (en cm)
% #6 = couleur ellipses et flèche
% #7 = contenu du nœud B (texte, formule, TikZ, etc.)
\newcommand{\FlecheEllipse}[7]{%
    \begin{tikzpicture}[overlay]
        % --- Nœud source A ---
        \node[draw,#6,ellipse,inner sep=2pt](#1){};

        % --- Nœud cible B ---
        \coordinate (B) at ($(#1) + (#4cm,#5cm)$);
        \node[draw,#6,ellipse,inner sep=2pt] (B-node) at (B) {#7};

        % --- Flèche de A vers B ---
        \draw[->,#6,thick] (#1) edge [out=#2,in=#3] (B-node);
    \end{tikzpicture}%
}